\section{USAJMO 2021}
        \begin{enumerate}
        	\item (\textit{USAJMO 2021}) $\textbf{Note:}$ For any geometry problem whose statement begins with an asterisk $(*)$, the first page of the solution must be a large, in-scale, clearly labeled diagram. Failure to meet this requirement will result in an automatic 1-point deduction.


 Let $\mathbb{N}$ denote the set of positive integers. Find all functions $f : \mathbb{N} \rightarrow \mathbb{N}$ such that for positive integers $a$ and $b,$
 \[f(a^2 + b^2) = f(a)f(b) \text{ and } f(a^2) = f(a)^2.\]


 

	\item (\textit{USAJMO 2021}) $\textbf{Note:}$ For any geometry problem whose statement begins with an asterisk $(*)$, the first page of the solution must be a large, in-scale, clearly labeled diagram. Failure to meet this requirement will result in an automatic 1-point deduction.


 Rectangles $BCC_1B_2,$ $CAA_1C_2,$ and $ABB_1A_2$ are erected outside an acute triangle $ABC.$ Suppose that
 \[\angle BC_1C+\angle  CA_1A+\angle AB_1B=180^{\circ}.\]Prove that lines $B_1C_2,$ $C_1A_2,$ and $A_1B_2$ are concurrent.


 

	\item (\textit{USAJMO 2021}) $\textbf{Note:}$ For any geometry problem whose statement begins with an asterisk $(*)$, the first page of the solution must be a large, in-scale, clearly labeled diagram. Failure to meet this requirement will result in an automatic 1-point deduction.


 An equilateral triangle $\Delta$ of side length $L>0$ is given. Suppose that $n$ equilateral triangles with side length 1 and with non-overlapping interiors are drawn inside $\Delta$, such that each unit equilateral triangle has sides parallel to $\Delta$, but with opposite orientation. (An example with $n=2$ is drawn below.)

                    \begin{figure}[H]
                        \centering
                        \includegraphics[width=0.5\textwidth]{Problem/USAJMO_I_2021_3_Problem_Figure_1.png}
                    \end{figure}
                    
Prove that
 \[n \leq \frac{2}{3} L^{2}.\]


 

	\item (\textit{USAJMO 2021}) $\textbf{Note:}$ For any geometry problem whose statement begins with an asterisk $(*)$, the first page of the solution must be a large, in-scale, clearly labeled diagram. Failure to meet this requirement will result in an automatic 1-point deduction.


 Carina has three pins, labeled $A, B$, and $C$, respectively, located at the origin of the coordinate plane. In a move, Carina may move a pin to an adjacent lattice point at distance $1$ away. What is the least number of moves that Carina can make in order for triangle $ABC$ to have area 2021? 


 (A lattice point is a point $(x, y)$ in the coordinate plane where $x$ and $y$ are both integers, not necessarily positive.) 


 

	\item (\textit{USAJMO 2021}) $\textbf{Note:}$ For any geometry problem whose statement begins with an asterisk $(*)$, the first page of the solution must be a large, in-scale, clearly labeled diagram. Failure to meet this requirement will result in an automatic 1-point deduction.


 A finite set $S$ of positive integers has the property that, for each $s \in S,$ and each positive integer divisor $d$ of $s$, there exists a unique element $t \in S$ satisfying $\text{gcd}(s, t) = d$. (The elements $s$ and $t$ could be equal.) Given this information, find all possible values for the number of elements of $S$.


 

	\item (\textit{USAJMO 2021}) $\textbf{Note:}$ For any geometry problem whose statement begins with an asterisk $(*)$, the first page of the solution must be a large, in-scale, clearly labeled diagram. Failure to meet this requirement will result in an automatic 1-point deduction.


 Let $n \geq 4$ be an integer. Find all positive real solutions to the following system of $2n$ equations: 
 \begin{align*} a_{1} &=\frac{1}{a_{2 n}}+\frac{1}{a_{2}}, & a_{2}&=a_{1}+a_{3}, \\ a_{3}&=\frac{1}{a_{2}}+\frac{1}{a_{4}}, & a_{4}&=a_{3}+a_{5}, \\ a_{5}&=\frac{1}{a_{4}}+\frac{1}{a_{6}}, & a_{6}&=a_{5}+a_{7}, \\ &\vdots \\ a_{2 n-1}&=\frac{1}{a_{2 n-2}}+\frac{1}{a_{2 n}}, & a_{2 n}&=a_{2 n-1}+a_{1} \end{align*}



 

\end{enumerate}